\chapter{Errors and Warnings}\label{chap:errors}

Every \hologo{LaTeX} user meets error messages, every day.
They look cryptic at first, but most of them come from a handful of causes,
and the message usually names the place.
This chapter shows how to read them and lists the common ones.

\section{Reading the log}\label{sec:log}

Everything \hologo{LaTeX} reports goes into the log file, \texttt{report.log} next to \texttt{report.tex}.
Overleaf lists the errors and warnings under \textbf{Logs and output files}, next to the Recompile button,
with a link to the line;
the full log is at the bottom of that list (\textbf{Raw logs}).
Editors such as TeXstudio and VS Code show them in a panel of their own.

An error starts with an exclamation mark:

\begin{codeblock}[style=shell]
! Undefined control sequence.
l.12 \sectoin
             {Introduction}
\end{codeblock}

\begin{itemize}
    \item The first line says what went wrong.
    \item \texttt{l.12} is the line number in the source file.
    \item The line is split where \hologo{LaTeX} noticed the problem:
        the last thing on the upper line, here \code|\sectoin|, is what it could not handle.
\end{itemize}

A few habits make errors much less painful:

\begin{itemize}
    \item \textbf{Fix the first error first.}
        One mistake often causes a series of follow-up errors that disappear with it.
    \item \textbf{Look a little before the line.}
        The line number says where \hologo{LaTeX} noticed the problem, which can be after the mistake,
        e.g.\ the end of the paragraph for a missing closing brace.
    \item \textbf{Compile often.}
        If it worked five minutes ago, the mistake is in what you changed since.
    \item \textbf{Narrow it down.}
        If you cannot find the cause, comment out half of the suspicious part with \code|%| and compile again,
        then half of the half, until the error is gone.
\end{itemize}

With several files (\autoref{sec:input}), the log names the file that was being read
(in parentheses: \texttt{(./sections/theory.tex}),
and Overleaf links to the right file.

\section{Common errors}\label{sec:common-errors}

An error stops the compilation.
In a terminal, \hologo{LaTeX} waits at a question mark;
type \texttt{x} and press Enter (\autoref{sec:compiling}).
Overleaf usually still shows a PDF, but everything after the error may be wrong in it.
The messages below are the ones you will meet most often, exactly as \hologo{LaTeX} writes them.

\subsection*{Undefined control sequence}

\begin{codeblock}[style=shell]
! Undefined control sequence.
l.4 \sectoin
            {Introduction}
\end{codeblock}

A command that \hologo{LaTeX} does not know.
Either its name is misspelled, as here,
or it comes from a package that is not loaded,
e.g.\ \code|\toprule| without \texttt{booktabs} or \code|\includegraphics| without \texttt{graphicx}.
To find out which package a command needs, search for the command on the web or in the package lists of \autoref{chap:packages-chapter}.

\subsection*{Missing \$ inserted}

\begin{codeblock}[style=shell]
! Missing $ inserted.
<inserted text>
                $
l.4 The raw data is in cooling_
\end{codeblock}

Something that only works in math mode was used outside of it:
\code|_| or \code|^| in the text (here in a file name, \autoref{sec:math-mode}),
or a math command such as \code|\alpha| without dollar signs.
The fix is either math mode, \code|$\alpha$|, or the escaped character, \code|\_|.
The same message appears for an empty line inside display math, which is not allowed.

\subsection*{\textbackslash begin\{...\} ended by \textbackslash end\{...\}}

\begin{codeblock}[style=shell]
! LaTeX Error: \begin{itemize} on input line 8 ended by \end{itemized}.
\end{codeblock}

An environment was closed with a different name, here a typo in the \code|\end|.
If the second name is \code|\end{document}|,
the environment was never closed at all:
\hologo{LaTeX} reached the end of the document with the \code|\begin| still open.
The line number of the \code|\begin| is in the message.

\subsection*{File ended while scanning use of \ldots}

\begin{codeblock}[style=shell]
! File ended while scanning use of \textit.
\end{codeblock}

A missing closing brace: the argument of \code|\textit| never ends.
Related messages are \texttt{Paragraph ended before \textbackslash textit was complete},
\texttt{Missing \} inserted} and \texttt{Extra \}, or forgotten \$},
all of them about braces (or dollar signs) that do not come in pairs.
A good editor highlights the matching brace when the cursor is on one.

\subsection*{File not found}

\begin{codeblock}[style=shell]
! LaTeX Error: File `boktabs.sty' not found.
\end{codeblock}

A package that does not exist (a typo, as here) or is not installed.
\hologo{TeX} Live and Overleaf have nearly all packages;
\hologo{MiKTeX} installs a missing one, if it is allowed to (\autoref{sec:miktex}).
For an image the message is a little different:

\begin{codeblock}[style=shell]
! Package pdftex.def Error: File `figures/cooling-curve.pdf' not found
\end{codeblock}

Check the folder, the spelling, the extension, and upper and lower case
(which matter on Linux, macOS and Overleaf, even when Windows does not care).

\subsection*{Extra alignment tab}

\begin{codeblock}[style=shell]
! Extra alignment tab has been changed to \cr.
\end{codeblock}

A row of a table has more cells (more \code|&|) than the column types in the argument of \texttt{tabular} allow.
The same happens in \texttt{align} with too many \code|&| in one line.

\subsection*{Double subscript}

\begin{codeblock}[style=shell]
! Double subscript.
l.4 $T_0_
         1$
\end{codeblock}

Two subscripts on one symbol, \code|T_0_1|.
Write the subscript as one group, \code|T_{0,1}|, or nest it, \code|T_{0_1}|.

\subsection*{Unicode character not set up}

\begin{codeblock}[style=shell]
! LaTeX Error: Unicode character (*$\leq$*) (U+2264)
               not set up for use with LaTeX.
\end{codeblock}

A character that \hologo{pdfLaTeX} has no letter for, usually a symbol pasted from another program or a PDF.
Use the command for it instead, here \code|$\leq$| (\autoref{sec:symbols}).

\subsection*{Missing \textbackslash end\{document\}}

\begin{codeblock}[style=shell]
! Emergency stop.
<*> report.tex
*** (job aborted, no legal \end found)
\end{codeblock}

The file ended without \code|\end{document}|.
It also happens when an \code|\end{document}| is commented out, or hidden inside an unclosed brace.

\subsection{Text pasted from elsewhere}\label{sec:pasted}

Text copied from a word processor, a website or a PDF brings typographic characters with it:
curly quotation marks (\texttt{“...”}), long dashes (\texttt{–} and \texttt{—}),
a minus sign that is not a hyphen, non-breaking spaces you cannot see.
In normal text, \hologo{LaTeX} handles quotation marks and dashes,
but in other places they cause trouble that is hard to see:

\begin{itemize}
    \item In code you copied, e.g.\ a command from a website:
        an en dash in place of a hyphen makes an option, a label or a file name a different one,
        e.g.\ \texttt{cooling–curves.pdf} is not \texttt{cooling-curves.pdf}, and the file is \enquote{not found}.
    \item In math mode, a pasted minus sign stops the compilation with
        \texttt{Unicode character $-$ (U+2212) not set up for use with LaTeX}.
        Type the minus as a hyphen, \code|$a - b$|.
    \item A straight \code|"| typed for quotation marks gives a closing mark at both ends (\autoref{sec:special-chars}).
\end{itemize}

When an error points at a line that looks perfectly fine, retype the suspicious characters by hand.

\section{Common warnings}\label{sec:common-warnings}

A warning does not stop the compilation; the PDF is complete.
But a warning often means that something in it is wrong,
so read them, at least before you hand a document in.

\subsection{Undefined references}\label{sec:undefined-refs}

A reference to a label that does not exist prints \texttt{??},
here because of a typo in the label:

\slidecode{04-ref-typo.tex}
\sliderender{04-ref-typo.pdf}
\codefile[style=shell]{\slidesdir/listings/04-ref-warning.txt}

A citation of a key that is not in the \texttt{.bib} file prints the key in bold:

\slidecode{04-cite-typo.tex}
\sliderender{04-cite-typo.pdf}
\codefile[style=shell]{\slidesdir/listings/04-cite-warning.txt}

Both warnings say \texttt{undefined} and name what is missing.
Take the name out of the warning and search for it, in the document and in the \texttt{.bib} file.
At the end, the log also says
\texttt{There were undefined references} or \texttt{Please (re)run Biber on the file}.

The other cause is that the document was not compiled often enough:
references need two runs, citations a run of \texttt{biber} in between (\autoref{sec:compiling}).
Overleaf and \texttt{latexmk} take care of that;
if the warning remains after they are done, it is a typo.

\subsection{Rerun to get cross-references right}

\begin{codeblock}[style=shell]
LaTeX Warning: Label(s) may have changed. Rerun to get cross-references right.
\end{codeblock}

Some numbers changed since the last run; compile once more.
\texttt{latexmk} and Overleaf do that by themselves.

\subsection{Label multiply defined}

\begin{codeblock}[style=shell]
LaTeX Warning: Label `fig:cooling' multiply defined.
\end{codeblock}

The same label is used twice, typically after copying a figure or a table.
References to it point at one of the two; give each a name of its own.

\subsection{Overfull and underfull boxes}

\begin{codeblock}[style=shell]
Overfull \hbox (12.74pt too wide) in paragraph at lines 32--34
\end{codeblock}

A line that \hologo{LaTeX} could not break well enough sticks out into the right margin, here by \qty{12.74}{pt}.
The usual causes are long words it cannot hyphenate:
file names, commands in typewriter font, URLs, or a word in a language \texttt{babel} does not know.
Rephrase the sentence, move the long word,
or tell \hologo{LaTeX} where it may hyphenate a word, \code|ther\-mo\-me\-ter|.
Overflows of less than a point are invisible and can be ignored.
The class option \texttt{draft} marks overfull lines with a black bar, so they are easy to find.

\begin{codeblock}[style=shell]
Underfull \hbox (badness 10000) in paragraph at lines 12--14
\end{codeblock}

A line with too much space in it.
The most common cause is \code|\\| at the end of a paragraph or twice in a row (\autoref{sec:spaces}).

\subsection{Other warnings}

\begin{description}
    \item[\texttt{Float too large for page}]
        A figure or table is taller than the page.
        Make it smaller, e.g.\ with \code|height=0.8\textheight|.
    \item[\texttt{Font shape ... undefined, using ... instead}]
        A font combination that does not exist, e.g.\ bold typewriter in some fonts.
        Usually harmless; check that the text looks as intended.
    \item[\texttt{Empty bibliography}]
        \texttt{biber} has not run yet, or nothing was cited.
\end{description}

\section{Getting help}\label{sec:help}

\begin{itemize}
    \item \textbf{Search the exact message}, without your own file and command names.
        Almost every beginner's error has been asked and answered on
        \myhref{https://tex.stackexchange.com}{tex.stackexchange.com}.
    \item \textbf{Read the manual of the package.}
        \code|texdoc <package>| opens it on your own installation,
        \texttt{ctan.org/pkg/<package>} on the web (\autoref{sec:packages}).
    \item \textbf{Make a minimal example.}
        Copy the document and delete everything that is not needed to produce the error:
        other chapters, packages, text.
        Often the cause becomes obvious on the way.
        If not, the minimal example (on TeX StackExchange called a \gls{mwe}) is what to post when you ask for help.
    \item \textbf{Ask someone}, with the error message and the minimal example.
    \item \textbf{AI chat tools} can explain an error message and suggest code, and are often right.
        But they also invent commands and package options that do not exist,
        and mix up old and new ways of doing things (\code|$$...$$|, \texttt{natbib}).
        Check what they suggest against the manual, and compile it before you trust it.
\end{itemize}
